\begin{afterword}
\summaryhead

The Catholic Church, perhaps the largest voluntary institution in the world, is held together by the fear of hell, by devotion to its ideas and leaders, and by conformity at weekly meetings. Self-described rationalists think themselves above such bonds, so for a simple relief effort the Pope could mobilize far more people than Richard Dawkins. While that holds, gains for atheism are partly hollow.

The author grants the Church's harms, then asks whether, counting only the good done, the average Catholic outdoes the average atheist by being more active. False beliefs may motivate more than true ones. But fear of hell is an imperfect motivator, and rationalists who cared a great deal, perhaps helped by techniques from CBT and meditation, might keep up; the author calls this a speculation.

The author suspects that the real source of the Church's output is people meeting in person, which is hard for a scattered community. The proposals are regular meetings for motivation, norms that applaud caring and expect useful work, and, on a guess that the Church spends over half its money on religion, a norm of giving 5\% of income to real causes.

\responsehead

The post asks more than it asserts, and it keeps most of its hedges. It states the obvious objection to its hope, that some false beliefs may motivate more than any true one, and concedes that it is possible. It labels its own proposal for matching the Church ``a further-future speculation.'' Its main hypothesis is marked ``I suspect,'' and its figure for Church spending ``I would venture.''

The premise behind the question, that believers do more good than unbelievers, is not argued in the post, but evidence then available supported it for religious practice in general. In an American survey of 2000, weekly attenders gave more than three times as much as people who rarely attended or had no religion, and were more likely to give and volunteer for nonreligious causes too. The post's main suspicion, that the work is done by meeting in person rather than by fear of hell, was later supported by Putnam and Campbell's survey research: people with many friends at church were likely to be more generous.

A few passing claims are stated more firmly than their support. The remark that ``theologians don't really believe that stuff'' has no hedge and no evidence; the post it links is about one conversation with a woman who was not a theologian. CBT and Zen meditation are both said to be ``experimentally shown to yield real improvements.'' That was true of CBT. For meditation, a review two years earlier had found mostly poor-quality studies and no firm conclusions. The link on running out of ``mental energy'' is to ego depletion, the theory that self-control draws on a limited store that runs down with use. That theory later failed large replications; the point about believers who sin does not depend on it.

One piece of information bears on the post's arithmetic. Its norm of 5\% for real causes is derived from the suggested 10\% tithe. Most American churchgoers give far less than a tithe, and among American denominations Catholics tend to give the least. The post does not claim otherwise, but a reader should know that the tithe is a target, not a measure of what believers give.

In short: a careful, mostly hedged question about whether believers do more good than unbelievers, candid about its own speculation, whose premise the survey evidence of its day supported.
\end{afterword}
